\chapter{Methodology} \label{chap:Methodology}
% This chapter presents the methodological choices that guided the development of this thesis. It explains the overall research approach, the iterative process used to build and refine the work, the scope boundaries that define what the project does and does not cover, and the criteria used to evaluate the final result.

\section{Research Approach}
This thesis follows a design-oriented and implementation-driven research approach, supported by iterative prototyping. Rather than treating the problem as a purely theoretical exercise, the work focuses on building and refining a practical framework for time-aligned music visualisation that responds to the needs identified in the literature review and in the reference tools analysed in the previous chapter. The research process is therefore grounded in making, testing, and adjusting concrete visual and interaction decisions as the project evolves.

The methodology combines conceptual analysis with software development. First, the problem space was studied through the comparison of existing tools, visualisation strategies, and alignment workflows. That understanding was then used to define a focused system design that could support layered, time-aligned representations of musical data. From that point onward, implementation and evaluation proceeded together: each prototype was used to verify whether the proposed interaction model, visual structure, and alignment logic were adequate for the thesis goals.

This approach is appropriate because the main contribution of the thesis is not only to describe a concept, but to demonstrate how such a concept can be materialised into a workable framework. As a result, the methodology emphasizes practical feasibility, visual clarity, and alignment consistency, while still remaining informed by the academic context established in the literature review.

\section{Iterative Development Process}
The development of LayerIt followed an iterative path made up of five prototype cycles before the final framework took shape. Rather than moving from a fixed specification directly to implementation, the project evolved through successive experiments in which each prototype was built to answer a specific question, evaluated for what it revealed, and then either discarded or used as the basis for the next step. Seen in this way, the value of the iterations is not only in the software they produced, but in the design knowledge they accumulated.

We started by creating \textbf{BeatMapJS}. Its goal was to explore a web-based beat-annotation interface and to test whether existing beat-tracking tools could be combined with a more interactive annotation workflow. The prototype was built in JavaScript and HTML using BeatThis and Wavesurfer.js, and it included an interactive spectrogram, region annotation, a tempo-annotation table, and import/export support for SV-compatible text files. The main lesson was that, although the system was functional, it did not contribute a meaningful new annotation capability beyond the integration of BeatThis. More importantly, it did not produce a reusable Python library, which, from the start, was an important requirement. The decision carried forward was that a reusable Python core would be mandatory.

The second cycle was \textbf{MEI-Basic-Edit}. Its goal was to test whether manual score alignment could support audio-to-score annotation by allowing score elements to be adjusted against an audio layer. The prototype was built in JavaScript and HTML with Verovio, MEI, and SVG, and it rendered a score below a static spectrogram while allowing users to drag score elements horizontally to align them with the audio representation. It also included an embedded XML editor for direct manipulation of the underlying score code. The key lesson was that interactive score editing was expensive both conceptually and technically: it increased UI complexity and required a level of implementation effort that was not justified by the results. At the same time, this cycle provided valuable familiarity with MEI, Verovio, and SVG structure. The decision carried forward was that SVG would remain a viable composition substrate, but interactive editing would not be the main direction.

The third cycle was \textbf{BeatPrecision}, a fork of Piano Precision. Its goal was to extend an existing GUI environment with beat-annotation features while preserving the score-alignment functionality already present in the software. The prototype was built in Qt/C++, integrating BeatThis and implementing horizontal score-alignment alongside the existing spectrogram view. While functional, the project revealed significant challenges. Piano Precision's (and by extension Sonic Visualizer's) codebase proved complex and difficult to manage. Long compile times and the sheer size of the codebase made each successive modification increasingly challenging to implement, manage, and debug.. More importantly, this iteration taught us a critical lesson: the need to distinguish between the interaction and visualization problems in BA. We realized we could not tackle both simultaneously and would need to choose a direction. Recognizing a clear gap in the field, we opted to focus on the visualization aspect, as we saw greater potential for the thesis to benefit the community in this area. 

The fourth cycle was 	extbf{Score-Alignment.py}. Its goal was to explore score-aligned spectrograms in a Python environment and to assess the Modusa framework as a possible basis for the project. The prototype was built in Python with Modusa and matplotlib and produced an interactive spectrogram. The lesson was that a truly fluid interactive visualization, with zooming, panning, a timeline, and playback control, was much harder to implement well than expected. Even though the prototype worked, the result was not satisfactory enough to justify continuing in that direction. The decision carried forward was to abandon real-time interactivity and instead prioritize static, high-quality, reproducible output.

The fifth cycle was 	extbf{Score-Alignment.js}. Its goal was to revisit the aligned-spectrogram idea using the lessons learned from the earlier MEI and beat-annotation prototypes. The prototype was built in JavaScript and HTML with Verovio, MEI, and Wavesurfer.js, and it displayed an MEI score alongside a spectrogram while storing beat timestamps in the MEI data and highlighting notes on interaction. The lesson from this final exploratory cycle was that the idea was promising, but it also confirmed the need for a tool that was strictly focused on visualization, written in Python, and organized in an object-oriented way so that it could be reused and extended more easily. The decision carried forward was that this combination of requirements defined the brief for LayerIt.

\begin{table}[H]
	\centering
	\caption{Summary of development iterations}
	\label{tab:development_iterations}
	\begin{tabular}{p{2.4cm} p{3.2cm} p{4.0cm} p{3.4cm} p{3.2cm}}
		\hline
		\textbf{Iteration} & \textbf{Goal} & \textbf{Built} & \textbf{Lesson} & \textbf{Decision carried forward} \\
		\hline
		BeatMapJS & Explore web-based beat annotation & JS/HTML, BeatThis, Wavesurfer.js, interactive spectrogram, region annotation, tempo table, SV-compatible import/export & No new annotation feature beyond BeatThis; no reusable Python library & Reusable Python core is mandatory \\
		MEI-Basic-Edit & Test manual score alignment against audio & JS/HTML, Verovio, MEI, SVG, drag-to-align score, XML editor & Interactive score editing is costly; gained MEI/Verovio/SVG familiarity & SVG as composition substrate; no interactive editing focus \\
		BeatPrecision & Extend GUI beat annotation in an existing score-aligned tool & Qt/C++, BeatThis, horizontal score follow-along & Codebase was hard to extend; interaction and visualization must be separated & Focus on visualization, not the interaction problem \\
		Score-Alignment.py & Explore score-aligned spectrograms in Python & Python, Modusa, matplotlib, interactive spectrogram & Fluid interactivity proved impractical & Prefer static, high-quality, reproducible output \\
		Score-Alignment.js & Revisit aligned spectrograms with richer score interaction & JS/HTML, Verovio, MEI, Wavesurfer.js, beat timestamps in MEI, note highlighting & Promising, but confirmed need for a Python OO visualization tool & Define the LayerIt brief \\
		\hline
	\end{tabular}
\end{table}

Taken together, these five cycles show a deliberate narrowing of scope. The project moved from an interactive beat-annotation GUI toward a reusable, object-oriented visualization framework. Each prototype clarified one part of the problem and ruled out another, so the final system could be built on a more stable understanding of what the thesis needed to solve.

\section{Scope Decision}
% This section should define the boundaries of the thesis. It should state clearly which features, visual layers, datasets, alignment modes, or interactions are included in the project, and which ones are excluded for reasons of feasibility, relevance, or time constraints. The goal is to make the contribution of the thesis precise and avoid implying that the work attempts to solve every possible problem in audio-to-score visualisation.

\section{Evaluation Criteria}
% This section should explain how the resulting work will be judged. It should describe the criteria used to assess whether the system is useful, correct, and aligned with the thesis goals, such as visual clarity, interaction quality, alignment consistency, flexibility, and implementation robustness. If there are tests, comparisons, or manual inspection steps planned, they should also be introduced here so the reader knows how the methodology will be validated.
